```latex
% ---------------------------------------------------------------------
% PHỤ LỤC
% ĐỀ TÀI: XÂY DỰNG HỆ THỐNG ĐỌC NHÃN MÁC VÀ CẢNH BÁO
% THÀNH PHẦN DỊ ỨNG THỰC PHẨM
% ---------------------------------------------------------------------

\chapter{Phụ lục}
\label{ch:phuluc}
\markboth{Phụ lục}{}

Phụ lục trình bày hướng dẫn cài đặt và chạy hệ thống Allergy Warning,
bảng tự kiểm trước khi nộp báo cáo và phần nhận xét của giảng viên
hướng dẫn. Các nội dung này hỗ trợ việc kiểm tra, trình diễn và
đánh giá sản phẩm.

% =====================================================================
% A. HUONG DAN CAI DAT VA CHAY CHUONG TRINH
% =====================================================================

\section*{A. Hướng dẫn cài đặt và chạy chương trình}
\phantomsection
\addcontentsline{toc}{section}{A. Hướng dẫn cài đặt và chạy chương trình}

Phần này trình bày các bước cơ bản để cài đặt và khởi chạy hệ thống
Allergy Warning trên máy tính Windows. Người thực hiện cần chuẩn bị
mã nguồn, Python, các thư viện phụ thuộc và cấu hình cơ sở dữ liệu
theo yêu cầu của dự án.

\subsection*{Bước 1. Chuẩn bị mã nguồn}

Tải hoặc sao chép mã nguồn của dự án về máy tính. Nếu mã nguồn đã
được lưu trên GitHub, có thể sử dụng lệnh sau sau khi thay đường dẫn
bằng địa chỉ kho mã nguồn thực tế:

\begin{lstlisting}[language=bash, caption={Tải mã nguồn dự án}]
git clone <URL_REPOSITORY>
cd allergy-warning-flask-main
\end{lstlisting}

Nếu đã có sẵn thư mục mã nguồn trên máy tính, có thể mở trực tiếp
thư mục đó bằng Visual Studio Code hoặc terminal.

\subsection*{Bước 2. Tạo môi trường Python}

Mở terminal tại thư mục dự án và tạo môi trường ảo để quản lý các
thư viện sử dụng trong ứng dụng:

\begin{lstlisting}[language=bash, caption={Tạo và kích hoạt môi trường ảo}]
python -m venv .venv
.venv\Scripts\activate
\end{lstlisting}

Sau khi kích hoạt thành công, cài đặt các thư viện được khai báo
trong tệp \code{requirements.txt} nếu dự án có tệp này:

\begin{lstlisting}[language=bash, caption={Cài đặt các thư viện phụ thuộc}]
python -m pip install -r requirements.txt
\end{lstlisting}

Nếu dự án chưa có tệp \code{requirements.txt}, cần cài đặt các thư
viện theo danh sách phụ thuộc thực tế của mã nguồn.

\subsection*{Bước 3. Cấu hình cơ sở dữ liệu và dịch vụ AI}

Trước khi chạy ứng dụng, cần kiểm tra cấu hình kết nối cơ sở dữ
liệu và các thông tin cần thiết để sử dụng dịch vụ AI.

Các thông tin cần kiểm tra gồm:
\begin{itemize}
    \item Thông tin kết nối cơ sở dữ liệu được cấu hình trong dự án.
    \item Cơ sở dữ liệu và các bảng cần thiết đã được tạo.
    \item Khóa API của dịch vụ AI được cấu hình nếu ứng dụng sử dụng
    dịch vụ này.
    \item Các biến môi trường cần thiết đã được khai báo đúng theo
    cấu hình của dự án.
\end{itemize}

Không nên đưa mật khẩu cơ sở dữ liệu hoặc khóa API thật vào mã
nguồn công khai. Những thông tin này cần được lưu trong tệp cấu
hình riêng hoặc biến môi trường theo cách triển khai của dự án.

\subsection*{Bước 4. Khởi chạy ứng dụng}

Nếu tệp khởi chạy chính của dự án là \code{app.py}, có thể chạy
ứng dụng bằng lệnh:

\begin{lstlisting}[language=bash, caption={Khởi chạy ứng dụng Flask}]
python app.py
\end{lstlisting}

Nếu dự án sử dụng một tệp khởi chạy khác, cần thay lệnh trên bằng
lệnh tương ứng với cấu trúc mã nguồn thực tế.

Sau khi ứng dụng khởi động thành công, mở trình duyệt và truy cập
địa chỉ cục bộ được hiển thị trong terminal. Với cấu hình mặc định
thường dùng của Flask, địa chỉ có thể là:

\begin{center}
    \url{http://127.0.0.1:5000}
\end{center}

Địa chỉ và cổng truy cập cần được điều chỉnh theo thông báo thực tế
khi chạy ứng dụng.

\subsection*{Bước 5. Kiểm tra các chức năng chính}

Sau khi khởi chạy hệ thống, tiến hành kiểm tra các chức năng đã
được triển khai:
\begin{enumerate}
    \item Mở trang đăng nhập và kiểm tra khả năng truy cập hệ thống.
    \item Thử đăng ký tài khoản nếu chức năng này được bật.
    \item Truy cập giao diện chính và kiểm tra các chức năng được cung cấp.
    \item Tải lên ảnh nhãn thực phẩm để kiểm tra quy trình phân tích.
    \item Đối chiếu kết quả nhận diện thành phần với nhãn thực phẩm gốc.
    \item Kiểm tra thông tin dị nguyên và nội dung cảnh báo được hiển thị.
    \item Kiểm tra hồ sơ dị ứng và lịch sử quét nếu các chức năng này
    đã được triển khai.
\end{enumerate}

Kết quả kiểm tra cần được ghi nhận theo tình trạng chạy thực tế.
Nếu xảy ra lỗi kết nối cơ sở dữ liệu, lỗi dịch vụ AI hoặc lỗi xử
lý ảnh, cần kiểm tra thông báo trên terminal và cấu hình liên quan.

% =====================================================================
% B. BANG TU KIEM TRUOC KHI NOP
% =====================================================================

\section*{B. Bảng tự kiểm trước khi nộp}
\phantomsection
\addcontentsline{toc}{section}{B. Bảng tự kiểm trước khi nộp}

Trước khi nộp báo cáo, nhóm cần kiểm tra lại mã nguồn, tài liệu,
các chức năng và kết quả thực tế của hệ thống. Bảng dưới đây dùng
để ghi nhận những nội dung đã hoàn thành.

\begin{tukiem}
    \item Đã kiểm tra khả năng cài đặt và chạy hệ thống trên môi
    trường phù hợp.

    \item Đã chạy thử các chức năng chính trước khi trình diễn
    sản phẩm.

    \item Đã kiểm tra quy trình tải ảnh, phân tích nhãn và hiển
    thị kết quả.

    \item Đã kiểm tra việc đối chiếu thành phần với dữ liệu dị
    nguyên và cách hiển thị cảnh báo.

    \item Đã kiểm tra xử lý các trường hợp ảnh không hợp lệ hoặc
    quá trình phân tích gặp lỗi.

    \item Đã kiểm tra kết nối cơ sở dữ liệu và các chức năng lưu
    trữ có trong hệ thống.

    \item Đã chuẩn bị mã nguồn và hướng dẫn cài đặt để người khác
    có thể thực hiện lại theo các bước phù hợp.

    \item Đã kiểm tra các hình, bảng, sơ đồ và nội dung trích dẫn
    trong báo cáo.

    \item Đã đối chiếu các mục tiêu ở Chương~\ref{ch:modau} với
    kết quả và kết luận tại Chương~\ref{ch:ketluan}.

    \item Các số liệu đánh giá trong báo cáo đều dựa trên kết quả
    kiểm thử thực tế; những nội dung chưa kiểm thử được trình bày
    rõ là hạn chế.

    \item Đã kiểm tra tài liệu tham khảo, nguồn dữ liệu và các
    thư viện hoặc dịch vụ bên ngoài được sử dụng.

    \item Đã rà soát nội dung báo cáo, loại bỏ các đoạn hướng dẫn
    mẫu và thông tin không thuộc đề tài Allergy Warning.

    \item Đã kiểm tra lại thông tin sinh viên, giảng viên hướng
    dẫn, tên đề tài và các thông tin trên trang bìa.
\end{tukiem}

\begin{luuy}
Bảng tự kiểm cần được đánh dấu dựa trên tình trạng thực tế của
đồ án. Không đánh dấu hoàn thành đối với chức năng chưa được
kiểm tra hoặc chưa triển khai. Trước khi nộp, cần tìm trong
toàn bộ báo cáo những đoạn hướng dẫn mẫu còn sót lại và thay
bằng nội dung phù hợp với đề tài.
\end{luuy}

% =====================================================================
% C. NHAN XET CUA GIANG VIEN HUONG DAN
% =====================================================================

\section*{C. Nhận xét của giảng viên hướng dẫn}
\phantomsection
\addcontentsline{toc}{section}{C. Nhận xét của giảng viên hướng dẫn}

\khungnhanxet[7.5cm]

\vspace{0.8em}

\noindent
\begin{minipage}[t]{0.5\textwidth}
    \color{daunavy}
    {\color{daumaroon}\bfseries Điểm số}

    \vspace{0.7em}
    \dongchamnhan{Bằng số:}

    \vspace{0.9em}
    \dongchamnhan{Bằng chữ:}
\end{minipage}
\hfill
\okyten{\dauGVHD}{Giảng viên hướng dẫn}
```